\clearpage
\seccionInforme{4. Descripción del proceso de desarrollo}
\phantomsection\addcontentsline{toc}{section}{4. Descripción del proceso de desarrollo}

\vspace{0.2cm}
\phantomsection\addcontentsline{toc}{subsection}{4.1 Entorno de trabajo}
\colabtitle{4.1 Entorno de trabajo}

\begin{table}[H]
\centering\small
\caption{Entorno de desarrollo y ejecución}
\renewcommand{\arraystretch}{1.25}
\begin{tabular}{|>{\raggedright\arraybackslash}p{\dimexpr 0.280\linewidth-2\tabcolsep-2\arrayrulewidth\relax}|>{\raggedright\arraybackslash}p{\dimexpr 0.720\linewidth-2\tabcolsep-2\arrayrulewidth\relax}|}
\hline
\rowcolor{colabpyblue} \color{white}\textbf{Elemento} & \color{white}\textbf{Detalle} \\ \hline
Sistema operativo & Windows 11 Pro \\ \hline
\rowcolor{colabpybg} Editor & Visual Studio Code \\ \hline
Node.js / npm & 22.14.0 / 10.9.2 \\ \hline
\rowcolor{colabpybg} Expo CLI & 58.1.13 (proyecto con Expo SDK 58) \\ \hline
Emulador & Android Studio · Pixel 7a · Android 16 (API 36) · 1080 × 2400 px, 420 dpi (411 × 914 dp) \\ \hline
\rowcolor{colabpybg} Cliente de ejecución & Expo Go 58.0.5 en el emulador \\ \hline
Control de versiones & Git y GitHub \\ \hline
\end{tabular}
\end{table}

\vspace{0.2cm}
\phantomsection\addcontentsline{toc}{subsection}{4.2 Creación del proyecto y estructura}
\colabtitle{4.2 Creación del proyecto y estructura}

El proyecto se creó con la plantilla en blanco de TypeScript de Expo y se agregaron las dependencias con \texttt{npx expo install}, que elige versiones compatibles con el SDK instalado:

\needspace{10\baselineskip}
\begin{lstlisting}[language=bash, title={Creación del proyecto e instalación de dependencias}]
npx create-expo-app@latest matrilab-rn --template blank-typescript
cd matrilab-rn
npx expo install @react-navigation/native @react-navigation/native-stack \
  @react-navigation/bottom-tabs react-native-screens react-native-safe-area-context \
  react-native-svg react-native-qrcode-svg expo-camera @expo/vector-icons
npx expo install react-native-web react-dom @expo/metro-runtime
npx expo install jest-expo jest @types/jest -- --save-dev
\end{lstlisting}
\par\penalty-100

El código se organizó por responsabilidad: dominio sin dependencias de React, estado, tema, componentes reutilizables, navegación y una carpeta de pantallas por rol.

\needspace{17\baselineskip}
\begin{lstlisting}[language=Texto, title={Estructura de carpetas del proyecto}]
matrilab-rn/
|-- App.tsx                  # proveedores: área segura, tema, estado y navegación
|-- app.json                 # configuración de Expo (nombre, orientación, permisos)
|-- src/
|   |-- domain/              # types.ts, reglas.ts (reglas puras), datos.ts (2026-B)
|   |-- state/AppStore.tsx   # Context + useReducer y selectores
|   |-- theme/               # tokens.ts y ThemeContext.tsx
|   |-- components/ui.tsx    # Tarjeta, Boton, Chip, Insignia, BarraProgreso, Kpi...
|   |-- navigation/          # RootNavigator.tsx y tipos.ts
|   `-- screens/             # LoginScreen, PerfilScreen y carpetas estudiante/,
|                            # docente/ y admin/ con las demás interfaces
|-- __tests__/reglas.test.ts # 31 pruebas unitarias
|-- herramientas/capturas-android.mjs  # recorrido automático de capturas
`-- informe/                 # este informe: LaTeX, DOCX, figuras y generadores
\end{lstlisting}
\par\penalty-100

\vspace{0.2cm}
\phantomsection\addcontentsline{toc}{subsection}{4.3 Arquitectura}
\colabtitle{4.3 Arquitectura}

La aplicación se dividió en capas (Figura 7). La capa de presentación solo conoce el estado a través de hooks; el estado aplica las reglas del dominio, y el dominio está formado por funciones puras que no importan nada de React ni de la plataforma. Esa separación permitió probar las reglas de negocio con Jest en un segundo, sin emulador.

\begin{figure}[H]
\centering
\safeimage{1.0\textwidth}{img/diagramas/arquitectura.png}{Arquitectura por capas de MatriLab EPIS}
\caption{Arquitectura por capas de MatriLab EPIS}
\label{fig:arquitectura}
\end{figure}

\vspace{0.2cm}
\phantomsection\addcontentsline{toc}{subsection}{4.4 Modelo de dominio y reglas de negocio}
\colabtitle{4.4 Modelo de dominio y reglas de negocio}

Los horarios se guardan en minutos desde la medianoche (por ejemplo, 07:00 es 420). Con esa representación la regla RN-02 de cruce de horario se reduce a comparar intervalos; los extremos no cuentan, de modo que un bloque que termina a las 08:40 no choca con otro que empieza a las 08:40.

\needspace{14\baselineskip}
\begin{lstlisting}[language=TSX, title={src/domain/reglas.ts · cruce de horario y semáforo}]
/** RN-02: dos bloques se cruzan si comparten día y sus intervalos se solapan (los extremos no cuentan). */
export function haySolape(a: Bloque, b: Bloque): boolean {
  return a.dia === b.dia && a.inicio < b.fin && b.inicio < a.fin;
}

/** Semáforo de vacantes: rojo sin cupo, ámbar con menos del 30 % libre, verde en otro caso. */
export function semaforo(t: Turno): Semaforo {
  const libres = cuposLibres(t);
  if (libres === 0) return 'rojo';
  return libres / t.cuposTotales < 0.3 ? 'ambar' : 'verde';
}
\end{lstlisting}
\par\penalty-100

\texttt{validar\allowbreak{}Matricula} aplica los controles en el mismo orden del proceso BPMN del LAB 01 y devuelve el primer motivo de rechazo con un mensaje legible: ventana cerrada, grupo ya matriculado, otro grupo del mismo curso, límite de laboratorios, cruce con un laboratorio, cruce con la teoría y grupo sin cupo. La reserva se hace con una función inmutable: si entre la consulta y la confirmación el cupo se agotó, no modifica nada.

\needspace{11\baselineskip}
\begin{lstlisting}[language=TSX, title={src/domain/reglas.ts · reserva atómica del cupo (RF-05)}]
export function reservarCupo(turnos: Turno[], turnoId: string): { turnos: Turno[]; reservado: boolean } {
  const objetivo = turnos.find((t) => t.id === turnoId);
  if (!objetivo || cuposLibres(objetivo) === 0) return { turnos, reservado: false };
  return {
    reservado: true,
    turnos: turnos.map((t) => (t.id === turnoId ? { ...t, cuposOcupados: t.cuposOcupados + 1 } : t)),
  };
}
\end{lstlisting}
\par\penalty-100

El pase QR contiene la cadena \texttt{MATRILAB|matrícula|grupo|estudiante|huella}. La huella es un FNV-1a de 32 bits sobre los cuatro primeros campos: si alguien altera el contenido, la huella deja de coincidir y el escáner rechaza el pase. No pretende ser un mecanismo criptográfico; en producción la verificación debe hacerse en el servidor.

\vspace{0.2cm}
\phantomsection\addcontentsline{toc}{subsection}{4.5 Estado global, tema y componentes propios}
\colabtitle{4.5 Estado global, tema y componentes propios}

El estado se maneja con \texttt{use\allowbreak{}Reducer} dentro de un contexto. El reductor es la única puerta para cambiar los cupos: la acción \texttt{MATRICULAR} reserva el cupo y registra la matrícula en la misma transición, lo que en la interfaz garantiza que dos toques seguidos no ocupen dos puestos.

\needspace{14\baselineskip}
\begin{lstlisting}[language=TSX, title={src/state/AppStore.tsx · acción MATRICULAR}]
case 'MATRICULAR': {
  if (!e.usuario) return e;
  const { turnos, reservado } = reservarCupo(e.turnos, a.turnoId);
  if (!reservado) return e;
  const id = `M-${String(e.matriculas.length + 1).padStart(4, '0')}`;
  return {
    ...e,
    turnos,
    matriculas: [...e.matriculas, { id, estudianteId: e.usuario.id, turnoId: a.turnoId, fecha: a.ahora, codigo: id, estado: 'confirmada' }],
  };
}
\end{lstlisting}
\par\penalty-100

Los componentes visuales propios se construyeron solo con primitivas del núcleo. El botón, por ejemplo, es un \texttt{Pressable} cuyo estilo depende de la variante y del estado \texttt{pressed}, con ripple nativo en Android y estados de accesibilidad:

\needspace{16\baselineskip}
\begin{lstlisting}[language=TSX, title={src/components/ui.tsx · componente Boton}]
<Pressable
  onPress={onPress}
  disabled={deshabilitado || cargando}
  accessibilityRole="button"
  accessibilityLabel={accessibilityLabel ?? titulo}
  accessibilityState={{ disabled: Boolean(deshabilitado), busy: Boolean(cargando) }}
  android_ripple={{ color: 'rgba(255,255,255,0.25)' }}
  style={({ pressed }) => [
    estilos.boton,
    { backgroundColor: fondo, borderColor: variante === 'fantasma' ? c.borde : fondo, opacity: deshabilitado ? 0.45 : pressed ? 0.85 : 1 },
    style,
  ]}
>
\end{lstlisting}
\par\penalty-100

\vspace{0.2cm}
\phantomsection\addcontentsline{toc}{subsection}{4.6 Desarrollo de las interfaces}
\colabtitle{4.6 Desarrollo de las interfaces}

A continuación se describe cada interfaz. En cada figura, el panel (a) es el mockup de la sección 2.4, el panel (b) es la captura real en el emulador Android y, cuando existe, el panel (c) muestra un estado alternativo de la misma pantalla.

\subsubtitulo{I1 · Inicio de sesión institucional}

\begin{figure}[H]
\centering
\safeimage{0.66\textwidth}{img/figuras/fig_i01.png}{Inicio de sesión: (a) mockup, (b) pantalla inicial, (c) validación del dominio institucional}
\caption{Inicio de sesión: (a) mockup, (b) pantalla inicial, (c) validación del dominio institucional}
\end{figure}

\textbf{Elementos.} Pantalla de entrada con la marca sobre el color primario y una tarjeta blanca centrada. Contiene el logotipo de la EPIS, los campos de correo y contraseña con iconos, el botón para mostrar u ocultar la contraseña, el interruptor \textit{Mantener la sesión iniciada}, el botón \textbf{Ingresar} y tres accesos de demostración por rol. Al ingresar un correo que no termina en \texttt{@unsa.\allowbreak{}edu.\allowbreak{}pe} aparece el mensaje de error del panel (c).

\textbf{Componentes y layout.} La pantalla combina \texttt{Safe\allowbreak{}Area\allowbreak{}View} (márgenes seguros), \texttt{Keyboard\allowbreak{}Avoiding\allowbreak{}View} y un \texttt{Scroll\allowbreak{}View} con \texttt{flex\allowbreak{}Grow: 1} y \texttt{justify\allowbreak{}Content: 'center'}, de modo que el formulario queda centrado pero sube cuando aparece el teclado. La tarjeta tiene \texttt{width: '100\%'} con \texttt{max\allowbreak{}Width: 440} para no estirarse en tabletas ni en la web. Los tres accesos de demostración son una fila (\texttt{flex\allowbreak{}Direction: 'row'}) de botones con \texttt{flex: 1}, que se reparten el ancho por igual. Como el fondo es oscuro, la pantalla fija \texttt{<Status\allowbreak{}Bar style=\allowbreak{}"light" /\allowbreak{}>}.

\needspace{12\baselineskip}
\begin{lstlisting}[language=TSX, title={LoginScreen.tsx · validación del formulario}]
const validar = () => {
  const limpio = correo.trim().toLowerCase();
  if (!/^[a-z0-9._-]+@unsa\.edu\.pe$/.test(limpio)) return setError('Usa tu correo institucional @unsa.edu.pe.');
  if (clave.length < 6) return setError('La contraseña debe tener al menos 6 caracteres.');
  const u = USUARIOS.find((x) => x.correo === limpio);
  if (!u || clave !== CLAVE_DEMO) return setError('Correo o contraseña incorrectos.');
  setError(null);
  entrar(u);
};
\end{lstlisting}
\par\penalty-100

\clearpage
\subsubtitulo{I2 · Panel del estudiante}

\begin{figure}[H]
\centering
\safeimage{0.66\textwidth}{img/figuras/fig_i02.png}{Panel del estudiante: (a) mockup, (b) tema claro, (c) tema oscuro después de matricularse en Plataformas Emergentes}
\caption{Panel del estudiante: (a) mockup, (b) tema claro, (c) tema oscuro después de matricularse en Plataformas Emergentes}
\end{figure}

\textbf{Elementos.} Muestra el saludo con avatar e iniciales, una tarjeta con el estado de la ventana de matrícula y su fecha de cierre, cuatro indicadores (laboratorios matriculados, créditos, vacantes libres en la oferta y laboratorios que aún puede tomar) y la lista \textit{Mis laboratorios}, donde cada matrícula tiene una franja con el color del curso y los botones \textbf{Ver pase QR} y \textbf{Anular}. En el panel (c) se ve la misma pantalla en tema oscuro, ya con dos laboratorios.

\textbf{Componentes y layout.} Un \texttt{Scroll\allowbreak{}View} con \texttt{Refresh\allowbreak{}Control} permite deslizar para actualizar. Los cuatro KPI están en una fila con \texttt{flex\allowbreak{}Wrap: 'wrap'} y cada \texttt{Kpi} usa \texttt{flex\allowbreak{}Basis: '47\%'} y \texttt{flex\allowbreak{}Grow: 1}: en un teléfono se forman dos columnas y en una pantalla ancha caben cuatro sin cambiar el código. La tarjeta de cada matrícula es una fila donde la franja de color es un \texttt{View} de 5 dp de ancho y el contenido ocupa el resto con \texttt{flex: 1}. Anular usa el \texttt{Alert} nativo con un botón de estilo destructivo.

\needspace{10\baselineskip}
\begin{lstlisting}[language=TSX, title={DashboardScreen.tsx · cuadrícula de indicadores}]
<View style={estilos.kpis}>
  <Kpi valor={String(activas.length)} etiqueta="Laboratorios matriculados" icono="flask-outline" />
  <Kpi valor={String(creditos)} etiqueta="Créditos con laboratorio" icono="ribbon-outline" color={c.marca} />
  ...
</View>
// kpis: { flexDirection: 'row', flexWrap: 'wrap', gap: espacio.m }
// kpi:  { flexBasis: '47%', flexGrow: 1 }
\end{lstlisting}
\par\penalty-100

\clearpage
\subsubtitulo{I3 · Catálogo de grupos con semáforo}

\begin{figure}[H]
\centering
\safeimage{0.66\textwidth}{img/figuras/fig_i03.png}{Catálogo: (a) mockup, (b) todos los grupos, (c) filtro Sin cruce de horario}
\caption{Catálogo: (a) mockup, (b) todos los grupos, (c) filtro Sin cruce de horario}
\end{figure}

\textbf{Elementos.} Lista todos los grupos de laboratorio agrupados por curso. Cada tarjeta muestra la letra del grupo, día y hora, ambiente y docente, la insignia del semáforo (\textit{Agotado}, \textit{Últimos cupos}, \textit{Disponible}), una barra de ocupación y el texto \textit{N de M PC libres}. Arriba hay un buscador y tres filtros. Con el filtro \textit{Sin cruce de horario} (panel c) desaparece el grupo C de Plataformas Emergentes, que coincide con la teoría de Gestión de Proyectos del jueves.

\textbf{Componentes y layout.} La lista es una \texttt{Section\allowbreak{}List} con una sección por curso; \texttt{render\allowbreak{}Section\allowbreak{}Header} dibuja el punto de color y el nombre del curso, y \texttt{sticky\allowbreak{}Section\allowbreak{}Headers\allowbreak{}Enabled} lo mantiene fijo mientras se desplazan sus grupos. Los filtros son chips dentro de un \texttt{Scroll\allowbreak{}View} horizontal, de modo que pueden agregarse más sin romper el diseño. En cada tarjeta, la letra del grupo es un cuadro de 44 × 44 dp centrado con \texttt{align\allowbreak{}Items} y \texttt{justify\allowbreak{}Content}, el texto toma el espacio libre con \texttt{flex: 1} y la insignia queda al final de la fila.

\needspace{13\baselineskip}
\begin{lstlisting}[language=TSX, title={CatalogoScreen.tsx · lista por secciones}]
<SectionList
  sections={secciones}                 // [{ curso, data: Turno[] }, ...]
  keyExtractor={(t) => t.id}
  renderItem={renderTurno}             // Tarjeta tocable con semáforo y barra
  stickySectionHeadersEnabled
  renderSectionHeader={({ section }) => (
    <View style={estilos.seccion}>{/* punto de color + nombre del curso */}</View>
  )}
  ListEmptyComponent={<Vacio icono="search-outline" titulo="Sin resultados" detalle="..." />}
/>
\end{lstlisting}
\par\penalty-100

\clearpage
\subsubtitulo{I4 · Detalle del grupo y confirmación}

\begin{figure}[H]
\centering
\safeimage{0.66\textwidth}{img/figuras/fig_i04.png}{Detalle del grupo: (a) mockup, (b) grupo B disponible, (c) grupo C bloqueado por cruce con la teoría (RN-02)}
\caption{Detalle del grupo: (a) mockup, (b) grupo B disponible, (c) grupo C bloqueado por cruce con la teoría (RN-02)}
\end{figure}

\textbf{Elementos.} Se abre como hoja modal desde el catálogo. Presenta el curso, las insignias de grupo y semáforo, una tarjeta con horario, ambiente, docente y equipamiento, y el \textbf{mapa de puestos}: una celda por computadora, en rojo las ocupadas y en verde las libres. Debajo, un aviso verde confirma que no hay cruces; si los hay, el aviso rojo explica con qué curso y en qué horario, y el botón \textbf{Confirmar matrícula} queda deshabilitado (panel c).

\textbf{Componentes y layout.} El mapa de puestos es un \texttt{View} con \texttt{flex\allowbreak{}Direction: 'row'}, \texttt{flex\allowbreak{}Wrap: 'wrap'} y \texttt{gap: 6}, cuyas celdas miden 46 × 40 dp. Con 20, 25 o 30 computadoras la cuadrícula se reacomoda sola en cualquier ancho. El botón de confirmación vive en un pie \textbf{fuera} del \texttt{Scroll\allowbreak{}View}, así siempre está visible. Al confirmar, un \texttt{Activity\allowbreak{}Indicator} simula la latencia del servidor y luego \texttt{navigation.\allowbreak{}replace} sustituye el detalle por el pase, para que Atrás no regrese a un formulario ya enviado.

\needspace{12\baselineskip}
\begin{lstlisting}[language=TSX, title={DetalleTurnoScreen.tsx · mapa de puestos}]
<View style={estilos.mapa}>   {/* flexDirection: 'row', flexWrap: 'wrap', gap: 6 */}
  {Array.from({ length: turno.cuposTotales }, (_, i) => (
    <View key={i} style={[estilos.pc, { backgroundColor: i < turno.cuposOcupados ? c.peligroSuave : c.exitoSuave }]}>
      <Ionicons name="desktop-outline" size={14} color={i < turno.cuposOcupados ? c.peligro : c.exito} />
      <Text style={estilos.numeroPc}>{i + 1}</Text>
    </View>
  ))}
</View>
// pc: { width: 46, height: 40, alignItems: 'center', justifyContent: 'center' }
\end{lstlisting}
\par\penalty-100

\clearpage
\subsubtitulo{I5 · Pase de laboratorio con QR}

\begin{figure}[H]
\centering
\safeimage{0.44\textwidth}{img/figuras/fig_i05.png}{Pase de laboratorio: (a) mockup, (b) pase emitido tras confirmar la matrícula}
\caption{Pase de laboratorio: (a) mockup, (b) pase emitido tras confirmar la matrícula}
\end{figure}

\textbf{Elementos.} Es la constancia digital que el estudiante muestra en la puerta del laboratorio. La parte superior, en color primario, lleva el curso, el grupo, el horario y el sello \textit{Confirmada}; la inferior, el código QR, el código de verificación de ocho caracteres y los datos del estudiante, el ambiente y la matrícula. El botón \textbf{Compartir pase} abre la hoja nativa de compartir.

\textbf{Componentes y layout.} La tarjeta del pase es un \texttt{View} con \texttt{overflow: 'hidden'} dividido en dos zonas. Entre ellas, la perforación se dibuja con dos círculos del color del fondo en \texttt{position: 'absolute'} (\texttt{left: -12} y \texttt{right: -12}) y una línea discontinua. El sello usa \texttt{position: 'absolute'} con \texttt{top} y \texttt{right}. El lado del QR se calcula con \texttt{use\allowbreak{}Window\allowbreak{}Dimensions} (\texttt{Math.\allowbreak{}min(220, width - 140)}) y los datos se ordenan en dos columnas con \texttt{flex\allowbreak{}Wrap} y \texttt{width: '50\%'} en cada dato.

\needspace{10\baselineskip}
\begin{lstlisting}[language=TSX, title={FichaQRScreen.tsx · perforación y código QR}]
<View style={estilos.perforacion}>
  <View style={[estilos.muesca, { left: -12, backgroundColor: c.fondo }]} />
  <View style={[estilos.guiones, { borderColor: c.borde }]} />
  <View style={[estilos.muesca, { right: -12, backgroundColor: c.fondo }]} />
</View>
<QRCode value={payloadQR(m)} size={ladoQR} color="#14202E" backgroundColor="#FFFFFF" />
// muesca: { position: 'absolute', width: 24, height: 24, borderRadius: 12 }
\end{lstlisting}
\par\penalty-100

\clearpage
\subsubtitulo{I6 · Horario semanal consolidado}

\begin{figure}[H]
\centering
\safeimage{0.66\textwidth}{img/figuras/fig_i06.png}{Horario: (a) mockup, (b) vista semanal, (c) vista de lista por día}
\caption{Horario: (a) mockup, (b) vista semanal, (c) vista de lista por día}
\end{figure}

\textbf{Elementos.} Consolida en una semana los laboratorios matriculados (bloques rellenos con el color del curso) y las clases de teoría (bloques con borde discontinuo). El selector \textit{Semana / Lista} cambia a una vista de tarjetas agrupadas por día, más cómoda para leer en un teléfono pequeño.

\textbf{Componentes y layout.} La cuadrícula se construye a mano porque React Native no tiene CSS Grid. Una columna de horas y cinco columnas de día se colocan en fila; el ancho de cada día sale de \texttt{use\allowbreak{}Window\allowbreak{}Dimensions} con un mínimo de 64 dp. Dentro de cada día, cada bloque usa \texttt{position: 'absolute'} con \texttt{top} igual a la hora de inicio y \texttt{height} igual a la duración, ambos convertidos con una constante de 58 dp por hora. Si el ancho mínimo no cabe, la cuadrícula se envuelve en un \texttt{Scroll\allowbreak{}View} horizontal.

\needspace{14\baselineskip}
\begin{lstlisting}[language=TSX, title={HorarioScreen.tsx · cuadrícula con posición absoluta}]
const disponible = width - espacio.l * 2 - CANAL_HORAS;
const anchoDia = Math.max(64, disponible / DIAS_VISIBLES.length);
const desborda = anchoDia * DIAS_VISIBLES.length > disponible + 1;
...
style={{
  top: (e.inicio - INICIO) / 60 * ALTO_HORA + 1,
  height: (e.fin - e.inicio) / 60 * ALTO_HORA - 2,
  borderStyle: lab ? 'solid' : 'dashed',
}}
...
{desborda ? <ScrollView horizontal>{cuadricula}</ScrollView> : cuadricula}
\end{lstlisting}
\par\penalty-100

\clearpage
\subsubtitulo{I7 · Panel del docente}

\begin{figure}[H]
\centering
\safeimage{0.44\textwidth}{img/figuras/fig_i07.png}{Panel del docente: (a) mockup, (b) panel con los grupos a cargo}
\caption{Panel del docente: (a) mockup, (b) panel con los grupos a cargo}
\end{figure}

\textbf{Elementos.} Resume la carga del docente: grupos a cargo, estudiantes inscritos sobre la capacidad, ocupación promedio y asistencias registradas en el día. Un botón destacado lleva al escáner y debajo se listan los grupos con semáforo, barra de ocupación y acceso a la \textbf{Nómina}.

\textbf{Componentes y layout.} Toda la pantalla es una sola \texttt{Flat\allowbreak{}List}. El saludo, los KPI y el botón van en \texttt{List\allowbreak{}Header\allowbreak{}Component}, de modo que se desplazan junto con la lista sin anidar una lista dentro de un \texttt{Scroll\allowbreak{}View}, situación que React Native desaconseja porque anula la virtualización. Las cabeceras de cada tarjeta usan \texttt{justify\allowbreak{}Content: 'space-between'}.

\needspace{10\baselineskip}
\begin{lstlisting}[language=TSX, title={PanelDocenteScreen.tsx · cabecera dentro de la lista}]
<FlatList
  data={mios}
  keyExtractor={(t) => t.id}
  ItemSeparatorComponent={() => <View style={{ height: espacio.m }} />}
  ListHeaderComponent={<View>{/* saludo, KPI y botón Escanear */}</View>}
  renderItem={({ item }) => <Tarjeta>{/* curso, semáforo, barra y botón Nómina */}</Tarjeta>}
/>
\end{lstlisting}
\par\penalty-100

\clearpage
\subsubtitulo{I8 · Nómina y asistencia}

\begin{figure}[H]
\centering
\safeimage{0.44\textwidth}{img/figuras/fig_i08.png}{Nómina del grupo: (a) mockup, (b) dos asistencias registradas con su hora}
\caption{Nómina del grupo: (a) mockup, (b) dos asistencias registradas con su hora}
\end{figure}

\textbf{Elementos.} Lista a los estudiantes del grupo ordenados alfabéticamente, con número, nombre y CUI. Tocar una fila registra la asistencia con la hora; la fila pasa de \textit{Pendiente} a una píldora verde. El resumen superior muestra presentes sobre el total, una barra de avance y un buscador por nombre o CUI. El pie ofrece \textbf{Exportar CSV} y \textbf{Escanear}.

\textbf{Componentes y layout.} Es una \texttt{Flat\allowbreak{}List} cuyo primer elemento, el resumen, queda fijo con \texttt{sticky\allowbreak{}Header\allowbreak{}Indices=\allowbreak{}\{[0]\}}. Las filas tienen alto constante (68 dp) para que el desplazamiento sea uniforme. Cada fila es un \texttt{Pressable} con \texttt{accessibility\allowbreak{}Role=\allowbreak{}"checkbox"} y \texttt{accessibility\allowbreak{}State} marcado, de modo que el lector de pantalla anuncia si el estudiante ya está presente. El título de la cabecera se cambia con \texttt{navigation.\allowbreak{}set\allowbreak{}Options}.

\needspace{15\baselineskip}
\begin{lstlisting}[language=TSX, title={NominaScreen.tsx · resumen fijo y filas marcables}]
<FlatList
  data={visibles}
  keyExtractor={(f) => f.id}
  stickyHeaderIndices={[0]}
  ListHeaderComponent={<View style={estilos.resumen}>{/* presentes, barra y buscador */}</View>}
  renderItem={({ item, index }) => (
    <Pressable onPress={() => marcar(item.id)} disabled={Boolean(estado.asistencias[item.id])}
      accessibilityRole="checkbox" accessibilityState={{ checked: Boolean(estado.asistencias[item.id]) }}>
      ...
    </Pressable>
  )}
/>
\end{lstlisting}
\par\penalty-100

\clearpage
\subsubtitulo{I9 · Escáner de asistencia}

\begin{figure}[H]
\centering
\safeimage{0.66\textwidth}{img/figuras/fig_i09.png}{Escáner: (a) mockup, (b) cámara activa en el emulador, (c) pase validado con el código de verificación}
\caption{Escáner: (a) mockup, (b) cámara activa en el emulador, (c) pase validado con el código de verificación}
\end{figure}

\textbf{Elementos.} Muestra la vista de la cámara con una guía de encuadre y lee el QR del pase. Si la cámara no está disponible, el docente puede digitar el código de verificación de ocho caracteres. El resultado aparece en una tarjeta verde con el nombre del estudiante, el curso y la hora, o en una roja si el pase es inválido o pertenece a otro docente.

\textbf{Componentes y layout.} El visor es un \texttt{View} de 340 dp de alto con \texttt{overflow: 'hidden'}. La \texttt{Camera\allowbreak{}View} lo cubre con \texttt{Style\allowbreak{}Sheet.\allowbreak{}absolute\allowbreak{}Fill} y encima se coloca otra capa con \texttt{absolute\allowbreak{}Fill} y \texttt{pointer\allowbreak{}Events=\allowbreak{}"none"}, que dibuja las cuatro esquinas del marco (cada una un \texttt{View} absoluto con dos bordes) y la leyenda. El campo manual normaliza lo escrito a mayúsculas hexadecimales y limita la longitud a ocho caracteres. Para evitar lecturas repetidas se bloquea el lector durante 2,5 s después de cada código.

\needspace{13\baselineskip}
\begin{lstlisting}[language=TSX, title={EscanerScreen.tsx · cámara con capa superpuesta}]
<View style={estilos.visor}>
  <CameraView style={StyleSheet.absoluteFill} facing="back"
    barcodeScannerSettings={{ barcodeTypes: ['qr'] }} onBarcodeScanned={alLeer} />
  <View style={[StyleSheet.absoluteFill, estilos.capa]} pointerEvents="none">
    <View style={estilos.marco}>
      {(['si', 'sd', 'ii', 'id'] as const).map((p) => <View key={p} style={[estilos.esquina, estilos[p]]} />)}
    </View>
    <Text style={estilos.ayuda}>Encuadra el código QR del pase</Text>
  </View>
</View>
\end{lstlisting}
\par\penalty-100

\clearpage
\subsubtitulo{I10 · Gestión de laboratorios}

\begin{figure}[H]
\centering
\safeimage{0.66\textwidth}{img/figuras/fig_i10.png}{Gestión de laboratorios: (a) mockup, (b) lista de ambientes, (c) edición en la hoja modal}
\caption{Gestión de laboratorios: (a) mockup, (b) lista de ambientes, (c) edición en la hoja modal}
\end{figure}

\textbf{Elementos.} La coordinación ve el total de computadoras operativas y una tarjeta por ambiente con ubicación, equipamiento, capacidad ajustable con botones − y +, barra de puestos reservados, interruptor de habilitación y botón \textbf{Editar datos}, que abre una hoja inferior para cambiar nombre, ubicación y equipamiento.

\textbf{Componentes y layout.} La lista es una \texttt{Flat\allowbreak{}List} con \texttt{num\allowbreak{}Columns} calculado a partir del ancho: una columna por debajo de 600 dp, dos hasta 960 dp y tres en adelante. \texttt{column\allowbreak{}Wrapper\allowbreak{}Style} agrega el espacio entre columnas. La edición usa el \texttt{Modal} nativo con fondo transparente: un \texttt{Pressable} semitransparente ocupa la pantalla (tocarlo cierra) y la hoja con esquinas superiores redondeadas queda abajo.

\needspace{14\baselineskip}
\begin{lstlisting}[language=TSX, title={LaboratoriosScreen.tsx · columnas según el ancho y edición modal}]
const { width } = useWindowDimensions();
const columnas = width < 600 ? 1 : width < 960 ? 2 : 3;

<FlatList
  key={columnas}
  data={estado.laboratorios}
  numColumns={columnas}
  columnWrapperStyle={columnas > 1 ? { gap: espacio.m } : undefined}
  ...
/>
<Modal visible={edicion !== null} animationType="slide" transparent onRequestClose={() => setEdicion(null)}>
\end{lstlisting}
\par\penalty-100

\clearpage
\subsubtitulo{I11 · Configuración de la matrícula}

\begin{figure}[H]
\centering
\safeimage{0.44\textwidth}{img/figuras/fig_i11.png}{Configuración de la matrícula: (a) mockup, (b) ventana abierta y reglas}
\caption{Configuración de la matrícula: (a) mockup, (b) ventana abierta y reglas}
\end{figure}

\textbf{Elementos.} Permite fijar la apertura y el cierre de la ventana de matrícula, el máximo de laboratorios por estudiante y dos reglas: permitir anular matrículas y avisar cuando un grupo se agota. Muestra el estado actual de la ventana y la lista de grupos agotados con la opción de solicitar un grupo adicional.

\textbf{Componentes y layout.} Es un formulario en un \texttt{Scroll\allowbreak{}View} con \texttt{keyboard\allowbreak{}Should\allowbreak{}Persist\allowbreak{}Taps=\allowbreak{}"handled"}, para que tocar un botón no solo cierre el teclado. Cada regla es una fila con \texttt{justify\allowbreak{}Content: 'space-between'}: texto a la izquierda con \texttt{flex: 1} y \texttt{Switch} a la derecha. Antes de guardar se valida el formato \texttt{dd/\allowbreak{}mm/\allowbreak{}aaaa hh:mm}, que el cierre sea posterior a la apertura y que el máximo esté entre 1 y 6; el error aparece sobre el botón.

\needspace{12\baselineskip}
\begin{lstlisting}[language=TSX, title={ConfiguracionScreen.tsx · validación de la ventana}]
const guardar = () => {
  const i = aISO(inicio);
  const f = aISO(fin);
  const max = Number(maximo);
  if (!i || !f) return setError('Usa el formato dd/mm/aaaa hh:mm en ambas fechas.');
  if (new Date(f) <= new Date(i)) return setError('El cierre debe ser posterior a la apertura.');
  if (!Number.isInteger(max) || max < 1 || max > 6) return setError('El máximo de laboratorios debe estar entre 1 y 6.');
  dispatch({ tipo: 'ACTUALIZAR_CONFIG', config: { inicio: i, fin: f, maxLaboratoriosPorEstudiante: max, ... } });
};
\end{lstlisting}
\par\penalty-100

\clearpage
\subsubtitulo{I12 · Perfil y ajustes}

\begin{figure}[H]
\centering
\safeimage{0.44\textwidth}{img/figuras/fig_i12.png}{Perfil y ajustes: (a) mockup, (b) perfil del estudiante}
\caption{Perfil y ajustes: (a) mockup, (b) perfil del estudiante}
\end{figure}

\textbf{Elementos.} Es común a los tres roles. Muestra los datos del usuario, el selector de tema (\textit{Sistema}, \textit{Claro}, \textit{Oscuro}), dos interruptores de notificaciones, información del dispositivo y el botón \textbf{Cerrar sesión}, que reinicia la pila de navegación hacia el inicio de sesión.

\textbf{Componentes y layout.} Las filas de ajustes reutilizan el componente \texttt{Fila\allowbreak{}Ajuste} (icono, texto con \texttt{flex: 1} y un control a la derecha). La sección \textit{Dispositivo} evidencia el entorno de ejecución: \texttt{Platform.\allowbreak{}OS} y \texttt{Platform.\allowbreak{}Version} indican el sistema, y \texttt{use\allowbreak{}Window\allowbreak{}Dimensions} el tamaño de la ventana en dp y la escala de texto. En el emulador se lee \textit{android · versión 36} y \textit{411 × 914 dp}. El tema elegido se guarda en el \texttt{Theme\allowbreak{}Context} y repinta toda la aplicación al instante.

\needspace{6\baselineskip}
\begin{lstlisting}[language=TSX, title={PerfilScreen.tsx · datos del dispositivo}]
<FilaAjuste icono="hardware-chip-outline" titulo="Plataforma" detalle={`${Platform.OS} · versión ${String(Platform.Version)}`} />
<FilaAjuste icono="resize-outline" titulo="Ventana"
  detalle={`${Math.round(width)} × ${Math.round(height)} dp · escala de texto ${fontScale.toFixed(2)}`} />
\end{lstlisting}
\par\penalty-100

\vspace{0.2cm}
\phantomsection\addcontentsline{toc}{subsection}{4.7 Adaptación a la orientación y a otras plataformas}
\colabtitle{4.7 Adaptación a la orientación y a otras plataformas}

Para comprobar que el layout responde al tamaño de la ventana se giró el emulador a horizontal. \texttt{use\allowbreak{}Window\allowbreak{}Dimensions} entrega entonces unos 915 dp de ancho y la lista de laboratorios pasa de una a dos columnas sin ningún código específico para la rotación (Figura 20). Para que esto fuera posible se cambió \texttt{orientation} a \texttt{"default"} en \texttt{app.\allowbreak{}json}, ya que con \texttt{"portrait"} el sistema bloqueaba el giro.

\begin{figure}[H]
\centering
\safeimage{0.85\textwidth}{img/capturas/13_laboratorios_horizontal.png}{Gestión de laboratorios en orientación horizontal: la FlatList pasa a dos columnas}
\caption{Gestión de laboratorios en orientación horizontal: la FlatList pasa a dos columnas}
\label{fig:horizontal}
\end{figure}

El mismo código se exportó a la web con \texttt{npx expo export --platform web}. \texttt{react-native-web} traduce los componentes al DOM y la aplicación funciona en el navegador (Figura 21); la versión publicada está en \url{https://owerlopez.github.io/matrilab-epis-react-native/}.

\begin{figure}[H]
\centering
\safeimage{0.98\textwidth}{img/figuras/fig_web.png}{La misma aplicación ejecutándose en un navegador web: inicio de sesión, panel del estudiante y catálogo}
\caption{La misma aplicación ejecutándose en un navegador web: inicio de sesión, panel del estudiante y catálogo}
\label{fig:web}
\end{figure}

\vspace{0.2cm}
\phantomsection\addcontentsline{toc}{subsection}{4.8 Pruebas automatizadas y verificación}
\colabtitle{4.8 Pruebas automatizadas y verificación}

Las reglas de negocio y el reductor se verificaron con \textbf{31 pruebas unitarias} en Jest [22], [23]. Se ejecutan con \texttt{npm test} y terminan en alrededor de un segundo porque no necesitan emulador. La Tabla 13 resume qué verifica cada grupo.

\begin{table}[H]
\centering\small
\caption{Pruebas unitarias agrupadas por regla verificada}
\renewcommand{\arraystretch}{1.25}
\begin{tabular}{|>{\raggedright\arraybackslash}p{\dimexpr 0.240\linewidth-2\tabcolsep-2\arrayrulewidth\relax}|>{\raggedright\arraybackslash}p{\dimexpr 0.100\linewidth-2\tabcolsep-2\arrayrulewidth\relax}|>{\raggedright\arraybackslash}p{\dimexpr 0.660\linewidth-2\tabcolsep-2\arrayrulewidth\relax}|}
\hline
\rowcolor{colabpyblue} \color{white}\textbf{Grupo} & \color{white}\textbf{Pruebas} & \color{white}\textbf{Qué se verifica} \\ \hline
Conversión de horas & 1 & \texttt{'07:00'} equivale a 420 minutos y 640 minutos a \texttt{'10:40'}. \\ \hline
\rowcolor{colabpybg} RN-02 cruce de horario & 3 & Solape el mismo día; bloques contiguos no se cruzan; días distintos no se cruzan. \\ \hline
Semáforo & 1 & Rojo sin cupo, ámbar con menos del 30 \% libre y verde en otro caso. \\ \hline
\rowcolor{colabpybg} validarMatricula & 8 & Acepta un grupo libre; rechaza cruce con teoría y con otro laboratorio, mismo grupo, mismo curso, grupo agotado, ventana cerrada y límite de laboratorios. \\ \hline
Reserva atómica (RF-05) & 4 & Suma un cupo sin mutar el arreglo; no reserva sin cupo; nunca supera el aforo con más solicitudes que vacantes; liberar no deja negativos. \\ \hline
\rowcolor{colabpybg} Pase QR & 3 & El contenido generado se verifica; un QR alterado se rechaza; la huella es estable y de 8 caracteres. \\ \hline
Horario & 1 & Los bloques se ordenan por día y luego por hora. \\ \hline
\rowcolor{colabpybg} Ventana y nómina & 3 & La ventana de demostración está abierta; la nómina tiene tantas filas como cupos; el CSV usa \texttt{;}. \\ \hline
Reductor & 7 & Matricular y anular actualizan cupos; sin sesión no se matricula; cerrar sesión; acciones de administración; código de verificación; la asistencia se marca una sola vez. \\ \hline
\end{tabular}
\end{table}

\needspace{16\baselineskip}
\begin{lstlisting}[language=bash, title={Resultado de npm run test:cobertura}]
PASS __tests__/reglas.test.ts
---------------|---------|----------|---------|---------|-------------------
File           | % Stmts | % Branch | % Funcs | % Lines | Uncovered Line #s
---------------|---------|----------|---------|---------|-------------------
All files      |    86.3 |    87.83 |   78.57 |   92.85 |
 domain        |   98.88 |    95.55 |     100 |     100 |
  datos.ts     |     100 |      100 |     100 |     100 |
  reglas.ts    |   98.52 |    95.34 |     100 |     100 | 132,140
 state         |   66.07 |    75.86 |      50 |   78.37 |
  AppStore.tsx |   66.07 |    75.86 |      50 |   78.37 | 76-84,95-99
---------------|---------|----------|---------|---------|-------------------
Tests:       31 passed, 31 total
Time:        1.088 s
\end{lstlisting}
\par\penalty-100

El dominio alcanza 100 \% de líneas cubiertas. En \texttt{App\allowbreak{}Store.\allowbreak{}tsx} quedan sin cubrir el proveedor de contexto y algunos selectores que solo se ejercitan al renderizar pantallas; esas rutas se verificaron en el recorrido manual y automatizado sobre Android descrito a continuación. Además, \texttt{npm run typecheck} (\texttt{tsc --no\allowbreak{}Emit}) termina sin errores en modo estricto.

\vspace{0.2cm}
\phantomsection\addcontentsline{toc}{subsection}{4.9 Ejecución en Android y recorrido automatizado de capturas}
\colabtitle{4.9 Ejecución en Android y recorrido automatizado de capturas}

La aplicación se ejecutó en el emulador con \texttt{npx expo start --android}, que inicia el servidor Metro y abre el proyecto en Expo Go. Para que las capturas del informe fueran reproducibles se escribió el script \texttt{herramientas/\allowbreak{}capturas-android.\allowbreak{}mjs}. El script recorre la aplicación por adb: con \texttt{uiautomator dump} obtiene el árbol de la pantalla, localiza los controles por su \texttt{accessibility\allowbreak{}Label} y los toca, escribe en los campos y guarda una captura de cada estado. Como depende de las etiquetas de accesibilidad y no de coordenadas, el mismo recorrido sirve en otro tamaño de pantalla y funciona además como prueba de humo de la navegación completa: iniciar sesión con cada rol, matricularse, ver el pase, marcar asistencia, validar el código en el escáner y girar la pantalla.

\needspace{15\baselineskip}
\begin{lstlisting}[language=TSX, title={herramientas/capturas-android.mjs · búsqueda de controles por accesibilidad}]
async function buscar(etiqueta, { exacto = true, intentos = 20, desplazar = false } = {}) {
  for (let i = 0; i < intentos; i++) {
    const lista = nodos();
    const coincide = (v) => (exacto ? v === etiqueta : v.includes(etiqueta));
    // Primero la etiqueta de accesibilidad (el control) y luego el texto visible (p. ej. un rótulo).
    const n = lista.find((x) => coincide(x.desc)) ?? lista.find((x) => coincide(x.texto));
    if (n && n.y > 0 && n.y < 2380) return n;
    if (desplazar) adb('shell', 'input', 'swipe', '540', '1700', '540', '800', '350');
    await dormir(700);
  }
  throw new Error(`No se encontró «${etiqueta}» en pantalla`);
}
\end{lstlisting}
\par\penalty-100

\vspace{0.2cm}
\phantomsection\addcontentsline{toc}{subsection}{4.10 Registro de cambios}
\colabtitle{4.10 Registro de cambios}

La guía pide dejar constancia de las modificaciones realizadas durante el trabajo. La Tabla 13 registra los cambios que surgieron al probar la aplicación en el emulador; en el documento DOCX cada cambio lleva además un comentario de revisión.

\begin{table}[H]
\centering\footnotesize
\caption{Cambios realizados durante el desarrollo y su motivo}
\renewcommand{\arraystretch}{1.25}
\begin{tabular}{|>{\raggedright\arraybackslash}p{\dimexpr 0.060\linewidth-2\tabcolsep-2\arrayrulewidth\relax}|>{\raggedright\arraybackslash}p{\dimexpr 0.220\linewidth-2\tabcolsep-2\arrayrulewidth\relax}|>{\raggedright\arraybackslash}p{\dimexpr 0.350\linewidth-2\tabcolsep-2\arrayrulewidth\relax}|>{\raggedright\arraybackslash}p{\dimexpr 0.370\linewidth-2\tabcolsep-2\arrayrulewidth\relax}|}
\hline
\rowcolor{colabpyblue} \color{white}\textbf{N.°} & \color{white}\textbf{Archivo} & \color{white}\textbf{Cambio} & \color{white}\textbf{Motivo} \\ \hline
1 & LoginScreen.tsx & Se agregó \texttt{<Status\allowbreak{}Bar style=\allowbreak{}"light" /\allowbreak{}>}. & La hora y los iconos del sistema quedaban oscuros sobre el fondo azul. \\ \hline
\rowcolor{colabpybg} 2 & LaboratoriosScreen.tsx & Una columna por debajo de 600 dp (antes dos desde 400 dp). & En el Pixel 7a (411 dp) las tarjetas quedaban comprimidas y el stepper se cortaba. \\ \hline
3 & app.json & \texttt{orientation} pasó de \texttt{"portrait"} a \texttt{"default"}. & El sistema bloqueaba el giro y no se podía comprobar el layout horizontal. \\ \hline
\rowcolor{colabpybg} 4 & EscanerScreen.tsx & El campo manual solo acepta hexadecimal en mayúsculas y 8 caracteres. & Evitar rechazos por minúsculas o espacios al dictar el código. \\ \hline
5 & NominaScreen.tsx & Se retiró \texttt{get\allowbreak{}Item\allowbreak{}Layout}. & La cabecera fija tiene alto variable y desplazaba los índices calculados. \\ \hline
\rowcolor{colabpybg} 6 & tsconfig.json & Se declararon los tipos \texttt{jest} y \texttt{node}. & TypeScript 6 no incluye tipos globales por defecto y \texttt{tsc} fallaba en las pruebas. \\ \hline
7 & reglas.test.ts & Se agregaron 8 pruebas (de 23 a 31). & Cubrir los motivos YA\_MATRICULADO y cruce con laboratorio, y las acciones del reductor. \\ \hline
\rowcolor{colabpybg} 8 & app.json & Se agregó \texttt{experiments.\allowbreak{}base\allowbreak{}Url} para la web. & Publicar la exportación web en GitHub Pages bajo una subruta. \\ \hline
\end{tabular}
\end{table}

\begin{notaInfo}
Cambio 2: la regla de columnas se ajustó después de ver las tarjetas cortadas en el emulador; el umbral de 600 dp coincide con el límite habitual entre teléfono y tableta pequeña.
\end{notaInfo}

\begin{notaInfo}
Cambio 3: con la orientación libre, la Figura 20 demuestra la adaptación automática de la rejilla.
\end{notaInfo}

\begin{notaOk}
Cambios 6 y 7: el comando \texttt{npm run typecheck} y las 31 pruebas se ejecutan sin errores después de estos ajustes.
\end{notaOk}

\vspace{0.2cm}
\phantomsection\addcontentsline{toc}{subsection}{4.11 Repositorio, ejecución y versión web}
\colabtitle{4.11 Repositorio, ejecución y versión web}

Todo el código, las capturas y este informe están en el repositorio \url{https://github.com/OwerLopez/matrilab-epis-react-native}. La versión web se publica automáticamente con GitHub Actions en \url{https://owerlopez.github.io/matrilab-epis-react-native/}. La carpeta \texttt{informe/\allowbreak{}} contiene las fuentes LaTeX de este documento, la versión DOCX, los diagramas en HTML y los scripts que los generan. Para ejecutar el proyecto:

\needspace{10\baselineskip}
\begin{lstlisting}[language=bash, title={Comandos de instalación y ejecución}]
git clone https://github.com/OwerLopez/matrilab-epis-react-native.git
cd matrilab-epis-react-native
npm install
npx expo start            # escanear el QR con Expo Go o pulsar «a» para el emulador Android
npm test                  # 31 pruebas unitarias
npm run typecheck         # verificación de tipos
npm run export:web        # versión web en dist/
\end{lstlisting}
\par\penalty-100

Credenciales de demostración: \texttt{olopeza@unsa.\allowbreak{}edu.\allowbreak{}pe}, \texttt{cmedina@unsa.\allowbreak{}edu.\allowbreak{}pe} o \texttt{laboratorios.\allowbreak{}epis@unsa.\allowbreak{}edu.\allowbreak{}pe} con la contraseña \texttt{unsa2026}, o los botones \textit{Estudiante}, \textit{Docente} y \textit{Admin} del inicio de sesión.

